HodgeCY II studies source-level fidelity for the frozen v1.0.0 double-octic
cohort.  Across 456 processed presentations, the census records 114
nontrivial source-level pairs or larger sets, split as 57 pairs, 13 triples,
and 44 larger sets.  The principal witness is the pair \(84/84a\): the two
presentations have the same local inventory, the same Hodge data
\((h^{1,1},h^{1,2})=(40,0)\), and the same rational source-assembly rank,
but their integral Smith types differ, namely
\((1^{23},2,6,12)\) and \((1^{21},2,4,4,4,12)\).
HodgeCY II then follows this source-level separation into a verified
block-geometry layer.  A self-contained eight-plane line-skeleton argument
identifies \(I_C=(F/L_1,\ldots,F/L_8)\), proves the Hilbert--Burch resolution
\[
0\to S(-8)^7\to S(-7)^8\to I_C\to 0,
\]
and uses regularity of the frozen quartic block section to compute the
degree-112 block scheme.  The proof allows higher-order point concurrences
while requiring no three planes to contain a common line.  For the block
schemes attached to \(84\) and \(84a\),
the degree-eight Hilbert/evaluation signatures agree: \(H_B(8)=105\) and
\(h^1({\bf P}^3,\mathcal I_B(8))=7\).  This structural collapse has Hilbert
series \((1-t^4)(1-8t^7+7t^8)/(1-t)^4\).  The result is a verified
non-determination theorem for this witness pair: degree-eight block-evaluation
data do not determine the integral source-assembly type.  The paper keeps the
ordinary-node promotion, classical nodal defect, source-to-evaluation morphism,
and LMHS/Hodge-atom interpretations open.
